\apendice{Pruebas y resultados}
\label{anexo:pruebas}

\section{Alcance y método de repetición}

Se compararon dos árboles Git de la aplicación:
\begin{description}
\item[C1] 19 de septiembre de 2026; \ruta{afb30bb6dddb584fabce86358ff679cd7a7e9aa7}.
\item[C2] 14 de octubre de 2026; \ruta{79d3a3de401b757f92c4ee2d6360628cbe17a37e}.
\end{description}
Las ejecuciones se efectuaron el 15 de octubre de 2026 en Windows, con Python 3.14, Node y Vitest para la interfaz, y Maven dentro de la imagen \ruta{maven:3.9.16-eclipse-temurin-26} para Java. Las pruebas de persistencia usaron contenedores locales reales de MySQL y PostgreSQL/pgvector. C1 se ejecutó en un worktree del commit indicado. C2 se ejecutó en el checkout \ruta{cd55d96}; se comprobó que su objeto Git \ruta{tree} coincide con el de \ruta{79d3a3d}. Así, los resultados C2 corresponden al mismo contenido, aunque no al mismo objeto de commit. No se modificó el producto para esta memoria.

Para repetir el ciclo, se obtiene un árbol limpio de cada SHA y se ejecutan las suites seleccionadas de la tabla~\ref{tabla:ciclos-prueba}. En Python se invoca \ruta{py -3.14 -m pytest -q} con los ficheros indicados; en la interfaz, \ruta{npx vitest run} con los ficheros indicados; en Java, \ruta{mvn test} con selección \ruta{-Dtest} para los módulos Core y worker, después de disponer de sus dependencias Maven. La integración de MySQL y PostgreSQL requiere Docker operativo. El nombre de cada fichero de prueba y el SHA hacen reproducible la selección; los XML de Surefire y la salida de pytest/Vitest son la evidencia de ejecución local, no una medición de producción.

\begin{table}[H]
\centering\small
\begin{tabularx}{\textwidth}{p{0.16\textwidth}p{0.33\textwidth}p{0.25\textwidth}X}
\toprule
\textbf{Corte} & \textbf{Selección} & \textbf{Resultado real} & \textbf{Alcance}\\
\midrule
C1 & Java: cuatro clases; scraper: tres ficheros; semántica: tres; web: tres. & 27 + 41 + 8 + 24 = 100 superadas; ninguna fallida. & Incluye una prueba MySQL y dos PostgreSQL.\\
C2 & Java: doce clases; scraper: cinco ficheros; semántica: cuatro; web: cinco. & 73 + 65 + 19 + 30 = 187 superadas; ninguna fallida. & Incluye dos pruebas MySQL y seis PostgreSQL.\\
\bottomrule
\end{tabularx}
\caption{Ciclos focalizados ejecutados el 15 de octubre de 2026.}
\label{tabla:ciclos-prueba}
\end{table}

En C1, las clases Java fueron \ruta{DownloadJobServiceTest}, \ruta{PublicHttpsUriPolicyTest}, \ruta{ZipArchiveBuilderTest} y \ruta{DownloadJobProcessorTest}. En C2 se añadieron \ruta{IdentityServiceTokenTest}, \ruta{IdentityControllerTest}, \ruta{CatalogControllerTest}, \ruta{SemanticSearchClientTest}, \ruta{DownloadDeliveryServiceTest}, \ruta{DownloadJobControllerTest}, \ruta{DownloadStorageBudgetTest} y \ruta{OperationalRetentionPrunerTest}. Core reside en \ruta{services/core-api}; el worker, en \ruta{services/download-worker}. Los resultados Java se comprobaron en los informes XML de Surefire de ambos módulos.

Para el scraper, C1 ejecutó \ruta{test\_validator.py}, \ruta{test\_safe\_http\_wrappers.py} y \ruta{test\_catalog\_projection\_mysql.py}; C2 añadió \ruta{test\_catalog\_projection.py} y \ruta{test\_description\_enricher.py}. Sus fuentes están en \ruta{api/scraper/tests}. Para semántica, C1 ejecutó \ruta{test\_model\_validation.py}, \ruta{test\_http\_policies.py} y \ruta{test\_store\_postgres.py}; C2 ejecutó \ruta{test\_search.py}, \ruta{test\_model\_validation.py}, \ruta{test\_indexer.py} y \ruta{test\_store\_postgres.py}, dentro de \ruta{services/semantic-service}. En la web, C1 ejecutó \ruta{AccountPages.test.tsx}, \ruta{AppSearchBar.test.tsx} y \ruta{i18n.test.tsx}; C2 añadió \ruta{account.test.ts} y \ruta{AdminScraperPage.test.tsx}. Estas pruebas están en el cliente de \ruta{services/webapp}. Los recuentos corresponden a estas selecciones, no a todas las pruebas del repositorio.

\section{Casos ligados a los riesgos}

Cada caso especifica condición, acción, expectativa y resultado observado. La referencia CP coincide con la tabla de trazabilidad~\ref{tabla:trazabilidad} y con el registro de riesgos~\ref{tabla:respuestas-riesgos}. \emph{Superado localmente} significa que pasó la selección indicada; \emph{parcial} o \emph{pendiente} impiden inferir el cierre del riesgo completo.

\begingroup\small\setlength{\tabcolsep}{3pt}
\begin{longtable}{p{0.08\textwidth}>{\raggedright\arraybackslash}p{0.42\textwidth}>{\raggedright\arraybackslash}p{0.40\textwidth}}
\caption{Casos de prueba y resultado observado en C2.}\label{tabla:casos-prueba}\\
\toprule
\textbf{ID} & \textbf{Precondición, acción y expectativa} & \textbf{Resultado, estado y evidencia}\\
\midrule
\endfirsthead
\toprule
\textbf{ID} & \textbf{Precondición, acción y expectativa} & \textbf{Resultado, estado y evidencia}\\
\midrule
\endhead
CP-01 & Dadas URL controladas con IP privada, HTTP o redirección insegura, validar el destino antes de descargar; debe rechazarse. & \textbf{Superado local.} Pasaron los validadores del scraper y \ruta{PublicHttpsUriPolicyTest}. No se atacó la web pública.\\
CP-02 & Dadas respuestas simuladas de proveedores, resolver y revalidar una fuente; una página o wrapper no debe aceptarse como instalador. & \textbf{Parcial.} Pasaron los validadores y la proyección del scraper. No se ensayó la disponibilidad futura de cada editor.\\
CP-03 & Dado un nombre o contenido no confiable, empaquetar solo entradas aceptadas y producir un manifiesto coherente; los casos parciales han de conservar estado. & \textbf{Superado local.} Pasaron \ruta{ZipArchiveBuilderTest} y \ruta{DownloadJobProcessorTest} en C1 y C2. No se inspeccionó cada instalador real.\\
CP-04 & Dados trabajos simultáneos y límites de espacio, aplicar presupuesto, cola y entrega autorizada; no liberar reserva antes de limpiar. & \textbf{Parcial.} Pasaron las pruebas de presupuesto, entrega y retención en C2. Faltan carga y fallo de disco.\\
CP-05 & Dado un enlace de acceso válido o ya usado y un trabajo de otro propietario, autenticar una sola vez y restringir operaciones. & \textbf{Parcial.} Pasaron las pruebas de identidad, acceso y trabajos en C2. No se certificó toda la autorización de extremo a extremo.\\
CP-06 & Dado MySQL real con historial de esquema, aplicar migraciones, cambiar fuentes y revisar proyección, contadores y rollback. & \textbf{Superado local.} Pasó la prueba MySQL del scraper: un caso en C1 y dos en C2. No se restauró una copia de producción.\\
CP-07 & Inventariar datos, terceros y retenciones y comprobar purga efectiva en sesiones, autenticación y auditoría. & \textbf{Pendiente en C2.} En ese corte no se acreditó una comprobación integral de la purga. La actualización posterior se recoge en la sección~\ref{sec:pruebas-retencion}.\\
CP-08 & Dado un índice ausente o incompleto, buscar y degradar la petición completa a literal; con índice sano, respetar filtros y página de MySQL. & \textbf{Superado local.} Pasaron búsqueda, indexador, PostgreSQL y clientes Core en C2. Incluye seis casos PostgreSQL reales; no mide latencia pública.\\
CP-09 & Dada una copia de los datos y volúmenes de producción, restaurar en un entorno aislado, verificar migraciones y abrir un trabajo completo. & \textbf{Pendiente.} Existen procedimiento y scripts, pero no evidencia de restauración completa ejecutada.\\
CP-10 & Dadas las dos revisiones, contrastar afirmaciones con código y pruebas, compilar fuentes y revisar páginas y referencias. & \textbf{Superado en documentación.} Los tres documentos compilan y se revisaron páginas clave; no se valida aquí el contrato OpenAPI.\\
CP-11 & Para cada editor, comparar licencia y términos de descarga con la automatización prevista antes de incorporarlo. & \textbf{Pendiente.} Se revisó el marco general, pero no se acredita autorización individual para todas las fuentes.\\
CP-12 & Dadas preferencias ES/EN y descripciones incompletas, cambiar idioma y detectar ausencia por idioma. & \textbf{Parcial.} Pasaron cuatro pruebas i18n en C1 y seis en C2, además de las de descripciones; faltan backfill completo y revisión humana.\\
\bottomrule
\end{longtable}
\endgroup

\section{Observación pública, incidencias y límites}

La consulta funcional de \url{https://batchdownloader.dev/} permitió observar la interfaz, el acceso por correo y las páginas \ruta{/privacy} y \ruta{/terms}.

No se registraron fallos en las selecciones automatizadas anteriores. En C2 quedaron abiertas como incidencias de verificación: la comprobación integral de la purga y revisión de terceros (R-07), la restauración completa (R-09), las condiciones por editor (R-11), el backfill bilingüe (R-12) y la capacidad bajo carga (R-03). Los resultados históricos del informe cronológico se utilizan para localizar cambios, pero no se suman a los recuentos ejecutados ni se presentan como comprobaciones actuales.

\section{Actualización de retención: 24 de octubre de 2026}
\label{sec:pruebas-retencion}

Esta selección posterior corresponde al árbol de trabajo de la rama \ruta{development}; se registra por separado de C1 y C2. Se ejecutaron pruebas Java con Maven, Temurin 25.0.5+8 y contenedores desechables, y pruebas Python de copias y Compose. Los recuentos de la tabla no duplican los casos repetidos durante la corrección.

\begin{table}[H]
\centering\small
\begin{tabularx}{\textwidth}{p{0.20\textwidth}p{0.14\textwidth}X}
\toprule
\textbf{Selección} & \textbf{Superadas} & \textbf{Comprobación}\\
\midrule
Core & 29 & Retención, descarte de correo caducado, retirada de destinatarios, baja de cuenta y coordinador de almacenamiento; incluye 22 casos con MySQL real.\\
Worker & 1 & Retirada de identificadores completados, conservando reservas asociadas a trabajos.\\
Notificaciones & 9 & Clasificación de errores sin persistir su contenido personal y limpieza de mensajes procesados o fallidos; incluye una prueba con H2 real.\\
Traducciones & 12 & Compilación y lectura de catálogos español e inglés y su contrato HTTP.\\
Copias & 7 & Fechas, límites y rutas; integra el script Bash con repositorios Restic locales y verifica que un fallo de listado impide ejecutar el borrado.\\
Compose & 18 & Agregación de salud, dependencias y capacidades de la composición.\\
\bottomrule
\end{tabularx}
\caption{Pruebas focalizadas de la política de retención: 51 Java y 25 Python.}
\end{table}
